% Options for packages loaded elsewhere
\PassOptionsToPackage{unicode}{hyperref}
\PassOptionsToPackage{hyphens}{url}
%
\documentclass[
  11pt,
]{article}
\usepackage{amsmath,amssymb}
\usepackage{iftex}
\ifPDFTeX
  \usepackage[T1]{fontenc}
  \usepackage[utf8]{inputenc}
  \usepackage{textcomp} % provide euro and other symbols
\else % if luatex or xetex
  \usepackage{unicode-math} % this also loads fontspec
  \defaultfontfeatures{Scale=MatchLowercase}
  \defaultfontfeatures[\rmfamily]{Ligatures=TeX,Scale=1}
\fi
\usepackage{lmodern}
\ifPDFTeX\else
  % xetex/luatex font selection
\fi
% Use upquote if available, for straight quotes in verbatim environments
\IfFileExists{upquote.sty}{\usepackage{upquote}}{}
\IfFileExists{microtype.sty}{% use microtype if available
  \usepackage[]{microtype}
  \UseMicrotypeSet[protrusion]{basicmath} % disable protrusion for tt fonts
}{}
\makeatletter
\@ifundefined{KOMAClassName}{% if non-KOMA class
  \IfFileExists{parskip.sty}{%
    \usepackage{parskip}
  }{% else
    \setlength{\parindent}{0pt}
    \setlength{\parskip}{6pt plus 2pt minus 1pt}}
}{% if KOMA class
  \KOMAoptions{parskip=half}}
\makeatother
\usepackage{xcolor}
\usepackage[margin=2.2cm]{geometry}
\usepackage{longtable,booktabs,array}
\usepackage{calc} % for calculating minipage widths
% Correct order of tables after \paragraph or \subparagraph
\usepackage{etoolbox}
\makeatletter
\patchcmd\longtable{\par}{\if@noskipsec\mbox{}\fi\par}{}{}
\makeatother
% Allow footnotes in longtable head/foot
\IfFileExists{footnotehyper.sty}{\usepackage{footnotehyper}}{\usepackage{footnote}}
\makesavenoteenv{longtable}
\setlength{\emergencystretch}{3em} % prevent overfull lines
\providecommand{\tightlist}{%
  \setlength{\itemsep}{0pt}\setlength{\parskip}{0pt}}
\setcounter{secnumdepth}{-\maxdimen} % remove section numbering
\ifLuaTeX
  \usepackage{selnolig}  % disable illegal ligatures
\fi
\IfFileExists{bookmark.sty}{\usepackage{bookmark}}{\usepackage{hyperref}}
\IfFileExists{xurl.sty}{\usepackage{xurl}}{} % add URL line breaks if available
\urlstyle{same}
\hypersetup{
  pdftitle={Worksheet 3D: Graphing Quadratics},
  hidelinks,
  pdfcreator={LaTeX via pandoc}}

\title{Worksheet 3D: Graphing Quadratics}
\author{}
\date{}

\begin{document}
\maketitle

\begin{quote}
\textbf{About This Worksheet}

This worksheet connects the algebra you've learned to visual
understanding. You'll see how changing numbers in an equation changes
the shape and position of a parabola.

\textbf{GeoGebra Links:} Throughout this worksheet, you'll find links to
interactive GeoGebra applets. Open these in a browser to explore the
concepts hands-on.
\end{quote}

\begin{center}\rule{0.5\linewidth}{0.5pt}\end{center}

\hypertarget{part-a-the-basic-parabola-8-problems}{%
\subsection{Part A: The Basic Parabola (8
problems)}\label{part-a-the-basic-parabola-8-problems}}

The simplest quadratic function is:

\[f(x) = x^2\]

** Explore:** \href{https://www.geogebra.org/calculator}{Basic Parabola}
\emph{Type \texttt{y\ =\ x\^{}2} into GeoGebra to see it.}

\begin{center}\rule{0.5\linewidth}{0.5pt}\end{center}

\textbf{1.} Complete this table for \(f(x) = x^2\):

\begin{longtable}[]{@{}llllllll@{}}
\toprule\noalign{}
\(x\) & -3 & -2 & -1 & 0 & 1 & 2 & 3 \\
\midrule\noalign{}
\endhead
\bottomrule\noalign{}
\endlastfoot
\(f(x)\) & & & & & & & \\
\end{longtable}

\textbf{2.} Plot these points on graph paper and connect them smoothly.

\textbf{3.} Answer these questions about \(f(x) = x^2\): - (a) What is
the lowest point on the graph? (This is called the \textbf{vertex}) -
(b) Is the parabola symmetric? Around what line? - (c) Does the parabola
open upward or downward? - (d) What is \(f(0)\)? This is the
\textbf{y-intercept}.

\textbf{4.} The graph of \(f(x) = x^2\) never goes below the x-axis.
What does this tell you about: - (a) The solutions to \(x^2 = 0\)? - (b)
The solutions to \(x^2 = -1\)?

\textbf{5.} Now consider \(g(x) = -x^2\). Complete the table:

\begin{longtable}[]{@{}llllllll@{}}
\toprule\noalign{}
\(x\) & -3 & -2 & -1 & 0 & 1 & 2 & 3 \\
\midrule\noalign{}
\endhead
\bottomrule\noalign{}
\endlastfoot
\(g(x)\) & & & & & & & \\
\end{longtable}

\textbf{6.} How does \(g(x) = -x^2\) compare to \(f(x) = x^2\)?

\textbf{7.} For \(g(x) = -x^2\): - (a) What is the vertex? - (b) Does it
open upward or downward? - (c) What is the maximum or minimum value?

\textbf{8.} \textbf{Rule:} If the coefficient of \(x^2\) is positive,
the parabola opens \_\_\_\_\_\_\_\_. If negative, it opens
\_\_\_\_\_\_\_\_.

\begin{center}\rule{0.5\linewidth}{0.5pt}\end{center}

\hypertarget{part-b-vertical-stretch-and-compression-8-problems}{%
\subsection{Part B: Vertical Stretch and Compression (8
problems)}\label{part-b-vertical-stretch-and-compression-8-problems}}

** Explore:** Open GeoGebra and graph these together: - \(y = x^2\)
(red) - \(y = 2x^2\) (blue)\\
- \(y = 0.5x^2\) (green) - \(y = 3x^2\) (orange)

\begin{center}\rule{0.5\linewidth}{0.5pt}\end{center}

\textbf{9.} Complete the table comparing \(f(x) = x^2\) and
\(g(x) = 2x^2\):

\begin{longtable}[]{@{}llllll@{}}
\toprule\noalign{}
\(x\) & -2 & -1 & 0 & 1 & 2 \\
\midrule\noalign{}
\endhead
\bottomrule\noalign{}
\endlastfoot
\(f(x) = x^2\) & & & & & \\
\(g(x) = 2x^2\) & & & & & \\
\end{longtable}

\textbf{10.} How does multiplying by 2 affect the y-values?

\textbf{11.} Compare \(f(x) = x^2\) and \(h(x) = 0.5x^2\):

\begin{longtable}[]{@{}llllll@{}}
\toprule\noalign{}
\(x\) & -2 & -1 & 0 & 1 & 2 \\
\midrule\noalign{}
\endhead
\bottomrule\noalign{}
\endlastfoot
\(f(x) = x^2\) & & & & & \\
\(h(x) = 0.5x^2\) & & & & & \\
\end{longtable}

\textbf{12.} How does multiplying by 0.5 affect the y-values?

\textbf{13.} Which parabola is ``narrower'' (steeper): \(y = x^2\) or
\(y = 3x^2\)?

\textbf{14.} Which parabola is ``wider'' (flatter): \(y = x^2\) or
\(y = 0.25x^2\)?

\textbf{15.} Match each equation to its description:

\begin{longtable}[]{@{}ll@{}}
\toprule\noalign{}
Equation & Description \\
\midrule\noalign{}
\endhead
\bottomrule\noalign{}
\endlastfoot
\(y = 5x^2\) & A. Same as \(y = x^2\) but flipped \\
\(y = 0.1x^2\) & B. Very narrow, opens up \\
\(y = -x^2\) & C. Very wide, opens up \\
\(y = -3x^2\) & D. Narrow, opens down \\
\end{longtable}

\textbf{16.} \textbf{Rule:} The coefficient \(a\) in \(y = ax^2\)
affects: - \textbf{Sign of \(a\):} Determines if parabola opens
\_\_\_\_\_ or \_\_\_\_\_ - \textbf{Size of \(|a|\):} Larger means
\_\_\_\_\_\_\_\_, smaller means \_\_\_\_\_\_\_\_

\begin{center}\rule{0.5\linewidth}{0.5pt}\end{center}

\hypertarget{part-c-vertical-shifts-8-problems}{%
\subsection{Part C: Vertical Shifts (8
problems)}\label{part-c-vertical-shifts-8-problems}}

** Explore:** Graph these in GeoGebra: - \(y = x^2\) (red) -
\(y = x^2 + 3\) (blue) - \(y = x^2 - 2\) (green)

\begin{center}\rule{0.5\linewidth}{0.5pt}\end{center}

\textbf{17.} Complete the table:

\begin{longtable}[]{@{}llllll@{}}
\toprule\noalign{}
\(x\) & -2 & -1 & 0 & 1 & 2 \\
\midrule\noalign{}
\endhead
\bottomrule\noalign{}
\endlastfoot
\(f(x) = x^2\) & & & & & \\
\(g(x) = x^2 + 3\) & & & & & \\
\(h(x) = x^2 - 2\) & & & & & \\
\end{longtable}

\textbf{18.} How does \(y = x^2 + 3\) compare to \(y = x^2\)?

\textbf{19.} How does \(y = x^2 - 2\) compare to \(y = x^2\)?

\textbf{20.} For each function, state the vertex: - (a) \(y = x^2\) →
Vertex: (\_\_, \_\_) - (b) \(y = x^2 + 5\) → Vertex: (\_\_, \_\_) - (c)
\(y = x^2 - 4\) → Vertex: (\_\_, \_\_) - (d) \(y = x^2 + 100\) → Vertex:
(\_\_, \_\_)

\textbf{21.} The equation \(y = x^2 + k\) shifts the parabola: - Up by
\(k\) units if \(k\) is \_\_\_\_\_\_\_\_ - Down by \(|k|\) units if
\(k\) is \_\_\_\_\_\_\_\_

\textbf{22.} Write the equation of a parabola that: - (a) Has the same
shape as \(y = x^2\) but vertex at \((0, 7)\) - (b) Has the same shape
as \(y = x^2\) but vertex at \((0, -3)\)

\textbf{23.} How many x-intercepts (roots) does each have? - (a)
\(y = x^2 + 1\) - (b) \(y = x^2\) - (c) \(y = x^2 - 4\)

\textbf{24.} Connect to algebra: Solve \(x^2 - 4 = 0\). How do the
solutions relate to the graph of \(y = x^2 - 4\)?

\begin{center}\rule{0.5\linewidth}{0.5pt}\end{center}

\hypertarget{part-d-horizontal-shifts-8-problems}{%
\subsection{Part D: Horizontal Shifts (8
problems)}\label{part-d-horizontal-shifts-8-problems}}

** Explore:** Graph these in GeoGebra: - \(y = x^2\) (red) -
\(y = (x - 2)^2\) (blue) - \(y = (x + 3)^2\) (green)

\begin{center}\rule{0.5\linewidth}{0.5pt}\end{center}

\textbf{25.} Complete the table for \(f(x) = (x - 2)^2\):

\begin{longtable}[]{@{}llllll@{}}
\toprule\noalign{}
\(x\) & 0 & 1 & 2 & 3 & 4 \\
\midrule\noalign{}
\endhead
\bottomrule\noalign{}
\endlastfoot
\((x-2)\) & & & & & \\
\(f(x) = (x-2)^2\) & & & & & \\
\end{longtable}

\textbf{26.} What is the vertex of \(y = (x - 2)^2\)?

\textbf{27.} Complete the table for \(g(x) = (x + 3)^2\):

\begin{longtable}[]{@{}llllll@{}}
\toprule\noalign{}
\(x\) & -5 & -4 & -3 & -2 & -1 \\
\midrule\noalign{}
\endhead
\bottomrule\noalign{}
\endlastfoot
\((x+3)\) & & & & & \\
\(g(x) = (x+3)^2\) & & & & & \\
\end{longtable}

\textbf{28.} What is the vertex of \(y = (x + 3)^2\)?

\textbf{29.} \textbf{The Tricky Part:} - \(y = (x - 2)^2\) shifts the
parabola \_\_\_\_\_\_ (left/right) by 2 units - \(y = (x + 3)^2\) shifts
the parabola \_\_\_\_\_\_ (left/right) by 3 units

\emph{The shift is OPPOSITE to the sign inside the parentheses!}

\textbf{30.} State the vertex of each: - (a) \(y = (x - 5)^2\) → Vertex:
(\_\_, \_\_) - (b) \(y = (x + 1)^2\) → Vertex: (\_\_, \_\_) - (c)
\(y = (x - 10)^2\) → Vertex: (\_\_, \_\_) - (d) \(y = (x + 7)^2\) →
Vertex: (\_\_, \_\_)

\textbf{31.} Write the equation of a parabola with: - (a) Vertex at
\((4, 0)\), same shape as \(y = x^2\) - (b) Vertex at \((-2, 0)\), same
shape as \(y = x^2\)

\textbf{32.} The equation \(y = (x - h)^2\) has vertex at \((h, 0)\).

If \(h = 3\), the vertex is at \_\_\_\_\_\_\_\_. If \(h = -4\), the
vertex is at \_\_\_\_\_\_\_\_.

\begin{center}\rule{0.5\linewidth}{0.5pt}\end{center}

\hypertarget{part-e-vertex-form-putting-it-all-together-10-problems}{%
\subsection{Part E: Vertex Form --- Putting It All Together (10
problems)}\label{part-e-vertex-form-putting-it-all-together-10-problems}}

The \textbf{vertex form} of a quadratic is:

\[y = a(x - h)^2 + k\]

Where: - \((h, k)\) is the \textbf{vertex} - \(a\) determines
\textbf{direction} (up/down) and \textbf{width}

** Explore:** \href{https://www.geogebra.org/m/BUD9HAPR}{Vertex Form
Explorer} \emph{Use sliders to change \(a\), \(h\), and \(k\) and watch
the parabola move!}

\begin{center}\rule{0.5\linewidth}{0.5pt}\end{center}

\textbf{33.} For \(y = 2(x - 1)^2 + 3\): - (a) What is \(a\)? \_\_\_ -
(b) What is \(h\)? \_\_\_ - (c) What is \(k\)? \_\_\_ - (d) What is the
vertex? (\_\_, \_\_) - (e) Does it open up or down? \_\_\_\_\_\_\_\_ -
(f) Is it narrower or wider than \(y = x^2\)? \_\_\_\_\_\_\_\_

\textbf{34.} For \(y = -0.5(x + 2)^2 - 1\): - (a) What is \(a\)? \_\_\_
- (b) What is \(h\)? \_\_\_ \emph{(Careful with the sign!)} - (c) What
is \(k\)? \_\_\_ - (d) What is the vertex? (\_\_, \_\_) - (e) Does it
open up or down? \_\_\_\_\_\_\_\_ - (f) Is it narrower or wider than
\(y = x^2\)? \_\_\_\_\_\_\_\_

\textbf{35.} Write the equation in vertex form for a parabola with: -
(a) Vertex \((3, 2)\), opens up, same width as \(y = x^2\) - (b) Vertex
\((-1, 4)\), opens down, same width as \(y = x^2\) - (c) Vertex
\((0, -5)\), opens up, twice as narrow as \(y = x^2\) - (d) Vertex
\((2, -3)\), opens down, half as steep as \(y = x^2\)

\textbf{36.} Match each equation to its vertex:

\begin{longtable}[]{@{}ll@{}}
\toprule\noalign{}
Equation & Vertex \\
\midrule\noalign{}
\endhead
\bottomrule\noalign{}
\endlastfoot
\(y = (x - 4)^2 + 1\) & A. \((4, 1)\) \\
\(y = (x + 4)^2 - 1\) & B. \((-4, -1)\) \\
\(y = (x - 4)^2 - 1\) & C. \((4, -1)\) \\
\(y = (x + 4)^2 + 1\) & D. \((-4, 1)\) \\
\end{longtable}

\textbf{37.} Sketch (without GeoGebra) the graph of
\(y = (x - 2)^2 - 4\): - (a) Vertex: \_\_\_\_\_\_\_\_ - (b) Opens:
\_\_\_\_\_\_\_\_ - (c) y-intercept (plug in \(x = 0\)): \_\_\_\_\_\_\_\_
- (d) x-intercepts (set \(y = 0\) and solve): \_\_\_\_\_\_\_\_

\textbf{38.} Sketch \(y = -(x + 1)^2 + 9\): - (a) Vertex:
\_\_\_\_\_\_\_\_ - (b) Opens: \_\_\_\_\_\_\_\_ - (c) y-intercept:
\_\_\_\_\_\_\_\_ - (d) x-intercepts: \_\_\_\_\_\_\_\_

\textbf{39.} A parabola has vertex \((2, 5)\) and passes through
\((0, 1)\). Find its equation.

\emph{Hint: Start with \(y = a(x - 2)^2 + 5\) and use the point to find
\(a\).}

\textbf{40.} A parabola has vertex \((-1, -4)\) and passes through
\((1, 0)\). Find its equation.

\textbf{41.} \textbf{Maximum/Minimum Values:}

For \(y = a(x - h)^2 + k\): - If \(a > 0\), the parabola has a
\_\_\_\_\_\_\_\_ value of \_\_\_ at \(x ={}\) \_\_\_ - If \(a < 0\), the
parabola has a \_\_\_\_\_\_\_\_ value of \_\_\_ at \(x ={}\) \_\_\_

\textbf{42.} A ball's height is \(h(t) = -5(t - 3)^2 + 50\). - (a) What
is the maximum height? - (b) When is the maximum height reached? - (c)
What was the initial height (\(t = 0\))?

\begin{center}\rule{0.5\linewidth}{0.5pt}\end{center}

\hypertarget{part-f-standard-form-and-converting-10-problems}{%
\subsection{Part F: Standard Form and Converting (10
problems)}\label{part-f-standard-form-and-converting-10-problems}}

The \textbf{standard form} is:

\[y = ax^2 + bx + c\]

To find the vertex from standard form:

\[x_{\text{vertex}} = \frac{-b}{2a}\]

Then plug this \(x\) back in to find \(y_{\text{vertex}}\).

\begin{center}\rule{0.5\linewidth}{0.5pt}\end{center}

\textbf{43.} For \(y = x^2 - 6x + 5\): - (a) Identify \(a\), \(b\),
\(c\): \_\_\_, \_\_\_, \_\_\_ - (b) Find
\(x_{\text{vertex}} = \frac{-b}{2a}\): \_\_\_ - (c) Find
\(y_{\text{vertex}}\) by plugging in: \_\_\_ - (d) Vertex: (\_\_, \_\_)

\textbf{44.} For \(y = x^2 + 4x + 3\): - (a) Find the vertex. - (b) Find
the y-intercept. - (c) Find the x-intercepts by factoring.

\textbf{45.} For \(y = -x^2 + 2x + 8\): - (a) Find the vertex. - (b)
Does this parabola have a maximum or minimum? - (c) Find the
x-intercepts.

\textbf{46.} For \(y = 2x^2 + 8x + 6\): - (a) Find the vertex. - (b)
Factor to find x-intercepts.

\textbf{47.} \textbf{Converting Standard to Vertex Form}

To convert \(y = x^2 - 4x + 7\) to vertex form, complete the square:

\(y = x^2 - 4x + 7\) \(y = (x^2 - 4x + 4) + 7 - 4\) ← Add and subtract
\((4/2)^2 = 4\) \(y = (x - 2)^2 + 3\)

Convert to vertex form: - (a) \(y = x^2 + 6x + 10\) - (b)
\(y = x^2 - 2x - 3\)

\textbf{48.} Convert to vertex form: - (a) \(y = x^2 + 8x + 20\) - (b)
\(y = x^2 - 10x + 21\)

\textbf{49.} \textbf{Converting Vertex to Standard Form}

Expand \(y = 2(x - 3)^2 + 1\):

\(y = 2(x^2 - 6x + 9) + 1 = 2x^2 - 12x + 18 + 1 = 2x^2 - 12x + 19\)

Convert to standard form: - (a) \(y = (x + 2)^2 - 5\) - (b)
\(y = 3(x - 1)^2 + 4\)

\textbf{50.} Convert to standard form: - (a) \(y = -(x - 4)^2 + 10\) -
(b) \(y = 2(x + 3)^2 - 7\)

\begin{center}\rule{0.5\linewidth}{0.5pt}\end{center}

\hypertarget{part-g-roots-and-the-discriminant-graphical-view-10-problems}{%
\subsection{Part G: Roots and the Discriminant --- Graphical View (10
problems)}\label{part-g-roots-and-the-discriminant-graphical-view-10-problems}}

The \textbf{x-intercepts} (roots) of a parabola are where \(y = 0\).

The \textbf{discriminant} \(b^2 - 4ac\) tells us how many roots to
expect:

\begin{longtable}[]{@{}
  >{\raggedright\arraybackslash}p{(\columnwidth - 4\tabcolsep) * \real{0.5000}}
  >{\raggedright\arraybackslash}p{(\columnwidth - 4\tabcolsep) * \real{0.2500}}
  >{\raggedright\arraybackslash}p{(\columnwidth - 4\tabcolsep) * \real{0.2500}}@{}}
\toprule\noalign{}
\begin{minipage}[b]{\linewidth}\raggedright
Discriminant
\end{minipage} & \begin{minipage}[b]{\linewidth}\raggedright
Roots
\end{minipage} & \begin{minipage}[b]{\linewidth}\raggedright
Graph
\end{minipage} \\
\midrule\noalign{}
\endhead
\bottomrule\noalign{}
\endlastfoot
\(b^2 - 4ac > 0\) & Two distinct real roots & Crosses x-axis twice \\
\(b^2 - 4ac = 0\) & One repeated root & Touches x-axis once (vertex on
x-axis) \\
\(b^2 - 4ac < 0\) & No real roots & Never crosses x-axis \\
\end{longtable}

** Explore:** \href{https://www.geogebra.org/m/mchqmhst}{Discriminant
Visualizer}

\begin{center}\rule{0.5\linewidth}{0.5pt}\end{center}

\textbf{51.} For each equation, calculate the discriminant and predict
the number of x-intercepts:

\begin{itemize}
\item
  \begin{enumerate}
  \def\labelenumi{(\alph{enumi})}
  \tightlist
  \item
    \(y = x^2 - 4x + 3\) → Discriminant: \_\_\_, x-intercepts: \_\_\_
  \end{enumerate}
\item
  \begin{enumerate}
  \def\labelenumi{(\alph{enumi})}
  \setcounter{enumi}{1}
  \tightlist
  \item
    \(y = x^2 - 4x + 4\) → Discriminant: \_\_\_, x-intercepts: \_\_\_
  \end{enumerate}
\item
  \begin{enumerate}
  \def\labelenumi{(\alph{enumi})}
  \setcounter{enumi}{2}
  \tightlist
  \item
    \(y = x^2 - 4x + 5\) → Discriminant: \_\_\_, x-intercepts: \_\_\_
  \end{enumerate}
\end{itemize}

\textbf{52.} Verify your predictions from problem 51 by: - (a) Factoring
or using quadratic formula for each - (b) Sketching each parabola (they
all have the same vertex x-coordinate!)

\textbf{53.} A parabola has equation \(y = x^2 + bx + 9\). For what
values of \(b\) does it: - (a) Cross the x-axis twice? - (b) Touch the
x-axis exactly once? - (c) Never touch the x-axis?

\textbf{54.} The graph of \(y = x^2 - 6x + k\) will touch the x-axis
exactly once when \(k ={}\) \_\_\_.

\textbf{55.} For \(y = ax^2 + 4x + 1\): - (a) If \(a = 1\), how many
x-intercepts? - (b) If \(a = 4\), how many x-intercepts? - (c) If
\(a = 5\), how many x-intercepts?

\textbf{56.} \textbf{Connecting Algebra and Geometry:}

The equation \(x^2 + 2x + 5 = 0\) has no real solutions.

Graphically, this means the parabola \(y = x^2 + 2x + 5\)
\_\_\_\_\_\_\_\_\_\_\_\_\_\_\_\_.

Find the vertex to confirm it's entirely above the x-axis.

\textbf{57.} Sketch a parabola that: - (a) Opens upward and has two
x-intercepts - (b) Opens downward and has two x-intercepts - (c) Opens
upward and has no x-intercepts - (d) Opens downward and has no
x-intercepts

\textbf{58.} A parabola has vertex \((3, -4)\) and \(a = 1\). - (a)
Write its equation. - (b) Will it cross the x-axis? Why? - (c) Find the
x-intercepts.

\textbf{59.} A parabola has vertex \((2, 5)\) and \(a = 1\). - (a) Write
its equation. - (b) Will it cross the x-axis? Why?

\textbf{60.} A parabola has vertex \((-1, 3)\) and \(a = -2\). - (a)
Write its equation. - (b) Will it cross the x-axis? Why? - (c) Find the
x-intercepts.

\begin{center}\rule{0.5\linewidth}{0.5pt}\end{center}

\hypertarget{part-h-real-world-graphing-applications-6-problems}{%
\subsection{Part H: Real-World Graphing Applications (6
problems)}\label{part-h-real-world-graphing-applications-6-problems}}

\textbf{61.} A ball is thrown with height \(h(t) = -5t^2 + 20t + 1\).

\begin{itemize}
\item
  \begin{enumerate}
  \def\labelenumi{(\alph{enumi})}
  \tightlist
  \item
    Write in vertex form.
  \end{enumerate}
\item
  \begin{enumerate}
  \def\labelenumi{(\alph{enumi})}
  \setcounter{enumi}{1}
  \tightlist
  \item
    What is the maximum height?
  \end{enumerate}
\item
  \begin{enumerate}
  \def\labelenumi{(\alph{enumi})}
  \setcounter{enumi}{2}
  \tightlist
  \item
    When is maximum height reached?
  \end{enumerate}
\item
  \begin{enumerate}
  \def\labelenumi{(\alph{enumi})}
  \setcounter{enumi}{3}
  \tightlist
  \item
    Sketch the graph for \(t \geq 0\).
  \end{enumerate}
\item
  \begin{enumerate}
  \def\labelenumi{(\alph{enumi})}
  \setcounter{enumi}{4}
  \tightlist
  \item
    When does the ball hit the ground?
  \end{enumerate}
\end{itemize}

\textbf{62.} A company's profit (in thousands) is modeled by
\(P(x) = -2x^2 + 24x - 54\) where \(x\) is units sold (in hundreds).

\begin{itemize}
\item
  \begin{enumerate}
  \def\labelenumi{(\alph{enumi})}
  \tightlist
  \item
    Find the vertex. What does it represent?
  \end{enumerate}
\item
  \begin{enumerate}
  \def\labelenumi{(\alph{enumi})}
  \setcounter{enumi}{1}
  \tightlist
  \item
    What sales level maximizes profit?
  \end{enumerate}
\item
  \begin{enumerate}
  \def\labelenumi{(\alph{enumi})}
  \setcounter{enumi}{2}
  \tightlist
  \item
    What is the maximum profit?
  \end{enumerate}
\item
  \begin{enumerate}
  \def\labelenumi{(\alph{enumi})}
  \setcounter{enumi}{3}
  \tightlist
  \item
    At what sales levels does the company break even (\(P = 0\))?
  \end{enumerate}
\end{itemize}

\textbf{63.} A farmer has 200m of fencing for a rectangular enclosure
against a barn (one side needs no fence). The area is:

\[A(w) = w(200 - 2w) = -2w^2 + 200w\]

\begin{itemize}
\item
  \begin{enumerate}
  \def\labelenumi{(\alph{enumi})}
  \tightlist
  \item
    What width maximizes the area?
  \end{enumerate}
\item
  \begin{enumerate}
  \def\labelenumi{(\alph{enumi})}
  \setcounter{enumi}{1}
  \tightlist
  \item
    What is the maximum area?
  \end{enumerate}
\item
  \begin{enumerate}
  \def\labelenumi{(\alph{enumi})}
  \setcounter{enumi}{2}
  \tightlist
  \item
    Sketch \(A(w)\) for sensible values of \(w\).
  \end{enumerate}
\end{itemize}

\textbf{64.} The demand for a product is \(D(p) = 1000 - 25p\) where
\(p\) is price in euros.

Revenue is price × quantity:
\(R(p) = p \cdot D(p) = p(1000 - 25p) = -25p^2 + 1000p\)

\begin{itemize}
\item
  \begin{enumerate}
  \def\labelenumi{(\alph{enumi})}
  \tightlist
  \item
    What price maximizes revenue?
  \end{enumerate}
\item
  \begin{enumerate}
  \def\labelenumi{(\alph{enumi})}
  \setcounter{enumi}{1}
  \tightlist
  \item
    What is the maximum revenue?
  \end{enumerate}
\item
  \begin{enumerate}
  \def\labelenumi{(\alph{enumi})}
  \setcounter{enumi}{2}
  \tightlist
  \item
    At what prices is revenue zero? Interpret these.
  \end{enumerate}
\end{itemize}

\textbf{65.} A suspension bridge cable forms a parabola. With the origin
at the lowest point, the cable follows \(y = 0.01x^2\) where \(x\) and
\(y\) are in meters.

\begin{itemize}
\item
  \begin{enumerate}
  \def\labelenumi{(\alph{enumi})}
  \tightlist
  \item
    How high is the cable at \(x = 50\)m from center?
  \end{enumerate}
\item
  \begin{enumerate}
  \def\labelenumi{(\alph{enumi})}
  \setcounter{enumi}{1}
  \tightlist
  \item
    If the towers are 100m apart, how high are they above the lowest
    point?
  \end{enumerate}
\item
  \begin{enumerate}
  \def\labelenumi{(\alph{enumi})}
  \setcounter{enumi}{2}
  \tightlist
  \item
    Sketch the cable.
  \end{enumerate}
\end{itemize}

\textbf{66.} Water from a fountain follows \(y = -0.5x^2 + 2x\) where
\(x\) is horizontal distance and \(y\) is height (meters).

\begin{itemize}
\item
  \begin{enumerate}
  \def\labelenumi{(\alph{enumi})}
  \tightlist
  \item
    What is the maximum height of the water?
  \end{enumerate}
\item
  \begin{enumerate}
  \def\labelenumi{(\alph{enumi})}
  \setcounter{enumi}{1}
  \tightlist
  \item
    How far horizontally does the water travel?
  \end{enumerate}
\item
  \begin{enumerate}
  \def\labelenumi{(\alph{enumi})}
  \setcounter{enumi}{2}
  \tightlist
  \item
    Sketch the water's path.
  \end{enumerate}
\end{itemize}

\begin{center}\rule{0.5\linewidth}{0.5pt}\end{center}

\hypertarget{geogebra-quick-reference}{%
\subsection{GeoGebra Quick Reference}\label{geogebra-quick-reference}}

\textbf{To graph in GeoGebra (geogebra.org/calculator):}

\begin{enumerate}
\def\labelenumi{\arabic{enumi}.}
\tightlist
\item
  Type equation directly: \texttt{y\ =\ x\^{}2} or
  \texttt{y\ =\ 2(x-3)\^{}2\ +\ 1}
\item
  Use \texttt{\^{}} for exponents: \texttt{x\^{}2} means \(x^2\)
\item
  Use \texttt{*} for multiplication: \texttt{2*x\^{}2} or just
  \texttt{2x\^{}2}
\item
  Add sliders: Type \texttt{a\ =\ 1} then use \(a\) in equations
\end{enumerate}

\textbf{Useful explorations:} - Type \texttt{y\ =\ a*x\^{}2} and create
slider for \(a\) - Type \texttt{y\ =\ (x\ -\ h)\^{}2\ +\ k} and create
sliders for \(h\) and \(k\) - Type \texttt{y\ =\ a*(x\ -\ h)\^{}2\ +\ k}
for full control

\end{document}
